\chapter{Conclusion}
\label{ch:con}

We formally proved most of the auxiliary lemmas Grobauer \cite{Grobauer} uses to show termination and correctness of the unification algorithm. We also discovered that formally proving some of the lemmas requires some more care than what is described in Grobauer's thesis. In particular, we showed that Lemma \ref{l22} and Facts \ref{f10} and \ref{f12} are not true for untyped lambda terms. We also showed that this extends to Nipkow's \cite{Nipkow-LICS-93} unification algorithm which does not always return the most general unifier for untyped terms. Furthermore, we implemented the unification algorithm in Isabelle and made a first effort towards proving its termination. We also discussed how the termination proof from Grobauer has to be adapted. Finally, we discussed advantages disadvantages of using de Bruijn indices for formal proofs.

Unfortunately, we were not able to formally prove the correctness of the unification algorithm. Completing the proof of correctness would be the obvious continuation of the work of this thesis. Moreover, one could try to formalize the enhanced versions of the unification algorithm which are detailed by Nipkow and contain efficiency improvements. Formalizing these extensions could be particularly interesting if the Isabelle specification was made executable. Then one could have an efficient, executable and verified unification algorithm for higher-order pattern unification.